\chapter{Conclusions}
\label{ch:conclusions}

\section{Summary}
\label{sec:conclusions_summary}

This thesis addressed how to interface a six-phase MP-VREH harvester with downstream power-management electronics, comparing three candidate rectifying topologies, the FWR star connection, the FWR delta connection, and the NVC, across the target $100$-$400\ \text{rpm}$ speed range. After introducing the harvester's operating principle and the candidate topologies (Chapters \ref{ch:theory}-\ref{ch:topologies}), each configuration was simulated (Chapter \ref{ch:simulations}), built and measured on a physical prototype board (Chapters \ref{ch:board_design}-\ref{ch:measurements}), and the simulated and measured results were compared (Chapter \ref{ch:comparison}) before evaluating compatibility with commercial PMICs (Chapter \ref{ch:pmic}).

\section{Key Findings}
\label{sec:conclusions_findings}

The NVC is the most efficient of the three configurations at every measured speed, while the FWR delta connection is consistently the least efficient and delivers the least power. The FWR star connection, however, overtakes the NVC in raw output power from around $250\ \text{rpm}$ onward, so the NVC and the star connection end up performing very similarly overall, each better suited to a different end of the speed range. The gaps between simulation and measurement are consistent with three causes: a real coil inductance somewhat lower than the nominal value used in simulation, the DC-bias derating of the ceramic compensation capacitors, and a loss mechanism that grows with electrical frequency and is not captured by the fixed-resistance coil model of Chapter \ref{ch:simulations}, as discussed in Section \ref{sec:comp_discussion}.

\section{Practical Implications}
\label{sec:conclusions_implications}

The delta connection and the NVC stay within the input voltage range of every  PMIC configuration across the whole speed range, while the star connection, despite giving the highest power at high speed, exceeds the input range of four of the six configurations above a certain speed and needs either a wide-range buck-boost PMIC or an input clamp to be used safely (Section \ref{sec:pmic_compatibility}). These compatibility figures are based on open-circuit voltage alone; verifying the total power conversion efficiency of a full harvester-PMIC chain was outside the scope of this work and is left as future work.

\section{Limitations}
\label{sec:conclusions_limitations}

The coil model used throughout the simulations represents losses with a single frequency-independent series resistance. In addition, the prototype's mechanical housing is 3D-printed, making it more sensitive to vibration and to small shifts in the alignment between the pickup units and the rotor than a machined structure would be; this makes highly repeatable, precise measurements difficult and likely contributes to part of the residual scatter between simulation and measurement.

\section{Future Work}
\label{sec:conclusions_future}

Promising directions include an adaptive compensation network that tracks the optimal $C_m$ across the whole speed range instead of a single fixed value, and testing the selected PMICs on hardware to measure the total power conversion efficiency of the complete chain from harvester to PMIC output, confirming whether the compatibility conclusions of Chapter \ref{ch:pmic} also hold in terms of delivered efficiency rather than voltage compatibility alone. Once this efficiency figure is established, the next step is to determine which downstream technology to implement and how to put this harvesting system to practical, useful application.
